\section{Resultados e métricas}

\subsection{Esquema de dados selecionado}
Para responder à pergunta analítica proposta com o menor consumo desnecessário de
memória, delimitamos o processamento estritamente a $7$ campos dos Microdados do
ENEM (Tabela~\ref{tab:esquema}).

\begin{table}[H]
  \centering
  \small
  \begin{tabularx}{\textwidth}{lXl}
    \toprule
    \textbf{Campo} & \textbf{Descrição} & \textbf{Papel no Processamento} \\
    \midrule
    \texttt{NU\_INSCRICAO} & Identificador do inscrito & Chave primária / controle de particionamento \\
    \texttt{SG\_UF\_PROVA} & Sigla da Unidade Federativa & Chave de agregação (27 estados) \\
    \texttt{TP\_ESCOLA} & Tipo de escola (1: Não inf., 2: Pública, 3: Privada) & Chave de agregação \\
    \texttt{TP\_PRESENCA\_MT} & Presença na prova de Matemática & Filtro booleano ($1 = \text{presente}$) \\
    \texttt{TP\_PRESENCA\_LC} & Presença na prova de Linguagens/Redação & Filtro booleano ($1 = \text{presente}$) \\
    \texttt{NU\_NOTA\_MT} & Nota da prova de Matemática ($0{,}0$ a $1000{,}0$) & Valor agregado ($\sum \text{nota}$, $\text{contagem}$) \\
    \texttt{NU\_NOTA\_REDACAO} & Nota da prova de Redação ($0{,}0$ a $1000{,}0$) & Valor agregado ($\sum \text{nota}$, $\text{contagem}$) \\
    \bottomrule
  \end{tabularx}
  \caption{Campos selecionados para o pipeline de dados.}
  \label{tab:esquema}
\end{table}

\subsection{Validação experimental do gerador e amostragem}
Executamos o utilitário \texttt{scripts/obter\_dados.py} para geração de uma
amostra de teste controlada de $10$\,MB com a semente determinística
$\text{seed}=42$. Os dados de execução e validação obtidos foram:

\begin{itemize}
  \item \textbf{Arquivo gerado:} \texttt{dados/enem\_sintetico\_amostra.csv}
  \item \textbf{Tamanho real em disco:} $10{,}07$\,MB ($10.562.878$ bytes)
  \item \textbf{Total de registros:} $325.000$ linhas
  \item \textbf{Candidatos presentes (Matemática):} $243.652$ ($\sim 74{,}97\%$)
  \item \textbf{Candidatos ausentes/eliminados:} $81.348$ ($\sim 25{,}03\%$)
  \item \textbf{Hash de integridade (SHA-256):}\\
    \texttt{\footnotesize 72ad34931db8e1adf99a53b0f99549d26a3c6075d98eafc7425ae76e0b5796b7}
\end{itemize}

\subsection{Taxa de redução no modelo MapReduce}
Um ponto crítico em Sistemas Distribuídos é a relação entre o volume de dados
transferido aos workers e o volume devolvido ao nó mestre. Com o esquema
definido:
\begin{itemize}
  \item Um bloco de dados (\emph{chunk}) típico de $100.000$ participantes
    possui cerca de $3{,}1$\,MB brutos.
  \item Após a filtragem e agregação local no worker, o resultado parcial
    consolidado consiste em uma tabela agregada de no máximo $27 \times 2 = 54$
    linhas (menos de $2$\,KB).
  \item Essa taxa de compressão lógica superior a $1.500:1$ garante que a rede
    não seja inundada na fase de agregação (\emph{gather}), tornando o ganho de
    paralelismo governado estritamente pela capacidade de computação dos nós.
\end{itemize}
